% Presentación de apertura, Congreso RIUSS, CIESS, 12 de noviembre de 2026.
% Base: presentación del 11 de noviembre de 2026 (versión final). Instrucción: BUILD-20261112-RIUSS-BEAMER v1.
% Archivo único: compila directo con pdflatex (dos pasadas); no necesita archivos externos.
% Clave de origen en cada lámina: C = base del 11 de noviembre, P = CEPAL (mayo 2026), K = arquitectura multipilar, N = nueva.
\documentclass[aspectratio=169,10pt]{beamer}
% Preámbulo: reconstrucción del estilo de la presentación del 22 de septiembre de 2026.
\usepackage[T1]{fontenc}
\usepackage[utf8]{inputenc}
\usepackage[scaled=0.92]{helvet}
\renewcommand{\familydefault}{\sfdefault}
\usepackage[spanish,es-nodecimaldot,es-noshorthands]{babel}
\usepackage{graphicx,booktabs,array,tabularx,amssymb}
\usepackage{tikz,pgfplots,pgfplotstable}
\pgfplotsset{compat=1.18}
\usetikzlibrary{positioning,arrows.meta,calc,patterns}

\definecolor{marino}{HTML}{18365C}
\definecolor{oxido}{HTML}{B0541C}
\definecolor{banda}{HTML}{EEF1F6}
\definecolor{barra}{HTML}{AEB9C6}
\definecolor{medio}{HTML}{5D728D}
\definecolor{pie}{HTML}{787878}
\definecolor{rejilla}{HTML}{E5E5E5}

\usetheme{default}
\setbeamersize{text margin left=1cm,text margin right=1cm}
\setbeamertemplate{navigation symbols}{}
\setbeamercolor{normal text}{fg=black}
\setbeamercolor{structure}{fg=marino}
\setbeamercolor{frametitle}{fg=marino}
\setbeamerfont{frametitle}{series=\bfseries,size=\large}
\setbeamertemplate{frametitle}{\vspace*{0.9em}\insertframetitle\par\vspace{0.75em}{\color{black}\hrule height 0.4pt}\vspace{0.2em}}
\setbeamertemplate{footline}{\hfill{\color{pie}\scriptsize\insertframenumber\,/\,\inserttotalframenumber}\hspace*{0.45cm}\vspace*{0.35cm}}
\setbeamertemplate{itemize item}{\color{marino}--}
\setbeamertemplate{itemize subitem}{\color{marino}--}
\setbeamertemplate{enumerate item}{\color{marino}\insertenumlabel.}
\setlength{\leftmargini}{1.3em}
\setlength{\leftmarginii}{1.3em}

\newcommand{\dest}[1]{{\color{oxido}\bfseries #1}}
\newcommand{\clave}[1]{{\color{marino}\bfseries #1}}
\newcommand{\fuente}[1]{\par\vspace{0.9em}{\color{pie}\scriptsize\raggedright #1\par}}
\newcommand{\notagraf}[1]{{\color{pie}\tiny\raggedright #1\par}}
\setbeamercolor{bandacolor}{bg=banda,fg=black}
\newcommand{\banda}[1]{%
  \par\vspace{0.7em}\noindent\makebox[\textwidth][c]{\begin{beamercolorbox}[wd=0.90\textwidth,sep=5pt,center]{bandacolor}\bfseries #1\end{beamercolorbox}}\par\vspace{0.7em}}
\setbeamercolor{cajacolor}{bg=banda,fg=black}
\newcommand{\caja}[2]{%
  \begin{beamercolorbox}[wd=\linewidth,sep=5pt]{cajacolor}
    {\color{marino}\bfseries #1}\par\vspace{0.25em}\footnotesize #2
  \end{beamercolorbox}}

\pgfplotsset{
  base/.style={axis lines=left,axis line style={black},tick style={black},
    tick label style={font=\tiny},label style={font=\tiny},
    ymajorgrids,grid style={rejilla},
    /pgf/number format/.cd,fixed,1000 sep={},/tikz/.cd,
    legend style={font=\tiny,draw=none,fill=none},legend cell align=left},
}

% Metadatos directos desde LaTeX: solo título y autor.
\hypersetup{pdftitle={La seguridad social en las Américas en el primer cuarto del siglo XXI},pdfauthor={Héctor Juan Villarreal Páez},pdfcreator={},pdfproducer={},pdfsubject={},pdfkeywords={}}
\pdfsuppressptexinfo=-1
\pdfinfoomitdate=1
\newcommand{\pef}{PEF 2027 (dictamen de la Cámara de Diputados)}

\title{La seguridad social en las Américas en el primer cuarto del siglo XXI}
\newcommand{\mx}{{\small\mdseries\color{oxido}Un ejemplo: México\enspace·\enspace}}

\begin{document}

% ===== 1. Portada (N) =====
\begin{frame}[plain]
  \centering
  \vspace*{0.9em}
  {\color{marino}\bfseries\Large La seguridad social en las Américas\\[0.15em] en el primer cuarto del siglo XXI\par}
  \vspace{0.4em}
  {\color{marino}\bfseries\normalsize Avances, transformaciones y desafíos persistentes\par}
  \vspace{1.7em}
  {\normalsize Héctor Juan Villarreal Páez\par}
  \vspace{0.9em}
  {\scriptsize ITED · Escuela de Gobierno y Transformación Pública · Tecnológico de Monterrey\\[0.3em]
  CONOSAFI · Comunidad Latinoamericana de Conocimiento en Sistemas de Salud y Financiamiento\par}
  \vspace{1.4em}
  {\small Congreso de la Red Internacional de Investigación en Seguridad Social (RIUSS)\\[0.2em]
  CIESS, Auditorio Benito Coquet, Ciudad de México · 12 de noviembre de 2026\par}
\end{frame}

% ===== 2. Tesis y ruta (N) =====
\begin{frame}{Tesis y ruta}
  \vspace{0.4em}
  \banda{En 25 años la región amplió la cobertura como nunca; no resolvió cómo pagarla durante la transición demográfica.}
  \vspace{0.1em}
  Voy a recorrer cinco temas:
  \vspace{0.3em}
  \begin{enumerate}\setlength{\itemsep}{0.3em}
    \item La \clave{demografía} de la región, que ya está definida para las próximas décadas.
    \item Las \clave{pensiones}: lo que avanzó, lo que cambió y lo que sigue pendiente.
    \item La \clave{salud} y la \clave{formalización}, con México como ejemplo del espacio fiscal.
    \item Una \clave{propuesta} de aseguramiento.
    \item \clave{CONOSAFI}, una comunidad regional para estudiar cómo se paga todo esto.
  \end{enumerate}
\end{frame}

% ===== 3. América Latina: menores de 15 y mayores de 65 (C3, regionalizada, con proyección propia) =====
\begin{frame}{América Latina: menores de 15 y mayores de 65}
  \vspace{0.6em}
  \begin{columns}[c,onlytextwidth]
    \begin{column}{0.44\textwidth}\small
      \begin{itemize}\setlength{\itemsep}{0.7em}
        \item En América Latina y el Caribe, la población de 65 años y más pasará de \dest{70 millones} en 2025 a \dest{138 millones} en 2050.
        \item Hacia \dest{2041} habrá más personas de 65 y más que menores de 15 años.
        \item Las personas de 80 y más casi se triplican: de \dest{13} a \dest{37 millones}. Es el grupo que más gasta en salud y en cuidados.
      \end{itemize}
    \end{column}
    \begin{column}{0.50\textwidth}\centering
      \begin{tikzpicture}
        \begin{axis}[base,width=\linewidth,height=4.6cm,xmin=2025,xmax=2051,ymin=0,ymax=170,
            xtick={2025,2030,2041,2050},ytick={0,50,100,150},ylabel={millones de personas},
            legend style={at={(0.5,1.02)},anchor=south,legend columns=2}]
          \draw[dashed,pie] (axis cs:2041,0) -- (axis cs:2041,170);
          \addplot[marino,thick] coordinates {(2025,144.3)(2030,131.5)(2035,120.7)(2040,112.5)(2045,110.0)(2050,106.1)};
          \addlegendentry{0 a 14 años}
          \addplot[oxido,thick] coordinates {(2025,69.9)(2030,83.1)(2035,96.2)(2040,109.6)(2045,123.4)(2050,137.8)};
          \addlegendentry{65 y más}
        \end{axis}
      \end{tikzpicture}
    \end{column}
  \end{columns}
  \fuente{Proyección propia para 49 países y territorios, con la fecundidad de cada país fija en su último valor observado. Población base y mortalidad: ONU, WPP 2024; México: INEGI, Encuesta Intercensal 2025. Detalle al final.}
\end{frame}

% ===== 4. La fecundidad cae más rápido de lo previsto (C4, regionalizada) =====
\begin{frame}{La fecundidad cae más rápido de lo previsto}
  \vspace{0.2em}
  \begin{columns}[c,onlytextwidth]
    \begin{column}{0.52\textwidth}\small
      \begin{itemize}\setlength{\itemsep}{0.45em}
        \item En los registros más recientes, Chile tiene \dest{0.99} hijos por mujer, Colombia \dest{1.01}, Costa Rica \dest{1.12}, y Argentina y México \dest{1.23}.
        \item Ponderada por población, la región está en alrededor de \dest{1.5} hijos por mujer. Para los mismos países, Naciones Unidas estima \dest{1.77} en 2025.
        \item Fernández-Villaverde ha mostrado que las proyecciones oficiales tienden a sobreestimar la fecundidad.
      \end{itemize}
      \vspace{0.4em}
      Por eso aquí usamos una proyección propia, con la fecundidad observada en cada país. Para 2050 da \dest{699 millones} de habitantes en la región; Naciones Unidas espera 730.
    \end{column}
    \begin{column}{0.45\textwidth}\centering
      \begin{tikzpicture}
        \begin{axis}[base,width=\linewidth,height=5.0cm,xbar=0.5pt,bar width=3.6pt,xmin=0,xmax=2.3,
            ymajorgrids=false,xmajorgrids,axis y line*=left,y axis line style={draw=none},
            ytick={1,...,8},ytick style={draw=none},
            yticklabels={Perú,Brasil,México,Argentina,Uruguay,Costa Rica,Colombia,Chile},
            ymin=0.4,ymax=8.6,xtick={0,0.5,1,1.5,2},xlabel={hijos por mujer},
            legend style={at={(0.5,1.02)},anchor=south,legend columns=2},area legend,reverse legend]
          \addplot[fill=barra,draw=marino] coordinates {(1.94,1)(1.60,2)(1.87,3)(1.50,4)(1.39,5)(1.31,6)(1.62,7)(1.13,8)};
          \addlegendentry{Naciones Unidas, 2025}
          \addplot[fill=oxido,draw=oxido] coordinates {(1.67,1)(1.52,2)(1.23,3)(1.23,4)(1.19,5)(1.12,6)(1.01,7)(0.99,8)};
          \addlegendentry{Observada}
        \end{axis}
      \end{tikzpicture}
    \end{column}
  \end{columns}
  \fuente{Observada: registros nacionales de 2022 a 2025 reunidos por Fernández-Villaverde (2026); México: INEGI, Encuesta Intercensal 2025. Naciones Unidas: WPP 2024, variante media. Fernández-Villaverde (2023).}
\end{frame}

% ===== 5. Por qué cayó tan rápido y por qué no esperamos un rebote (N; trabajo propio sobre el colapso de la fecundidad) =====
\begin{frame}{Por qué cayó tan rápido, y por qué no esperamos un rebote}
  \vspace{-0.5em}
  \begin{columns}[c,onlytextwidth]
    \begin{column}{0.58\textwidth}\small
      \begin{itemize}\setlength{\itemsep}{0.25em}
        \item En Costa Rica la fecundidad pasó de \dest{1.83} hijos por mujer en 2010 a \dest{1.12} en 2024. En Colombia, de \dest{1.7} a \dest{1.1} en menos de una década.
        \item No es un efecto de las cuentas: los nacimientos cayeron entre \dest{18} y \dest{47\%} mientras crecía el número de mujeres en edad fértil.
        \item Lo que cambió fue la composición de las parejas. En Colombia casi no cambió la proporción de mujeres de 20 a 39 años que vive en pareja, pero las casadas pasaron de \dest{20.7} a \dest{11.9\%}.
        \item No encontramos un contagio dentro de cada generación, sino \clave{reemplazo de generaciones}: cada una llega con menor propensión a casarse, y la conserva.
      \end{itemize}
      \vspace{0.1em}
      Una caída así es persistente. Por eso mantenemos la fecundidad observada.
    \end{column}
    \begin{column}{0.39\textwidth}\centering
      \begin{tikzpicture}
        \begin{axis}[base,width=\linewidth,height=4.0cm,ybar,bar width=20pt,ymin=0,ymax=26,ymajorgrids=false,
            ytick={0,5,10,15,20,25},ylabel={\% de las mujeres de 20 a 39 años},xmin=0.4,xmax=2.6,xtick={1,2},xtick style={draw=none},
            xticklabels={Al inicio de la serie,Al final de la serie}]
          \addplot[fill=barra,draw=marino,bar shift=0pt] coordinates {(1,20.7)};
          \addplot[fill=oxido,draw=oxido,bar shift=0pt] coordinates {(2,11.9)};
          \node[font=\tiny,anchor=south] at (axis cs:1,20.7) {20.7};
          \node[font=\tiny\bfseries,text=oxido,anchor=south] at (axis cs:2,11.9) {11.9};
        \end{axis}
      \end{tikzpicture}\par
      \notagraf{\centering Colombia: mujeres casadas. Encuesta de hogares, 2008--2024.}
    \end{column}
  \end{columns}
  \vspace{-0.8em}
  \fuente{Villarreal (2026), «Cohort Replacement, Not Cascade», documento de trabajo, con registros vitales de INEC y DANE y encuestas de hogares del DANE.}
\end{frame}

% ===== 5. La ventana demográfica se cierra (C5, regionalizada, con proyección propia) =====
\begin{frame}{La ventana demográfica se cierra}
  \vspace{0.5em}
  \begin{columns}[c,onlytextwidth]
    \begin{column}{0.46\textwidth}\small
      \begin{itemize}\setlength{\itemsep}{0.6em}
        \item En la región, la población de 15 a 64 años es hoy \dest{67.9\%} del total y llega a su máximo, \dest{68.5\%}, hacia \dest{2035}. Ese último aumento viene de que hay menos niños, no de que haya más trabajadores.
        \item Después baja: \dest{65.1\%} en 2050.
        \item Los países no van al mismo paso. Brasil ya va de bajada, Chile llega a su máximo hacia 2030 y Colombia hacia 2040. En 2050, Brasil y Chile estarán en \dest{63\%}.
      \end{itemize}
    \end{column}
    \begin{column}{0.50\textwidth}\centering
      \begin{tikzpicture}
        \begin{axis}[base,width=\linewidth,height=4.7cm,xmin=2025,xmax=2051,ymin=60,ymax=73,
            xtick={2025,2030,2035,2040,2045,2050},ytick={60,62,64,66,68,70,72},
            ylabel={\% de la población total},
            legend style={at={(0.5,1.02)},anchor=south,legend columns=4}]
          \addplot[oxido,very thick] coordinates {(2025,67.9)(2030,68.4)(2035,68.5)(2040,68.1)(2045,66.6)(2050,65.1)};
          \addlegendentry{Región}
          \addplot[marino,thick] coordinates {(2025,69.1)(2030,68.5)(2035,67.9)(2040,66.9)(2045,65.1)(2050,63.0)};
          \addlegendentry{Brasil}
          \addplot[marino,thick,dashed] coordinates {(2025,68.9)(2030,69.1)(2035,68.8)(2040,67.4)(2045,65.4)(2050,63.0)};
          \addlegendentry{Chile}
          \addplot[medio,thick,densely dotted] coordinates {(2025,69.8)(2030,70.4)(2035,70.9)(2040,71.6)(2045,69.7)(2050,67.4)};
          \addlegendentry{Colombia}
        \end{axis}
      \end{tikzpicture}
    \end{column}
  \end{columns}
  \fuente{Proyección propia para 49 países y territorios, con la fecundidad de cada país fija en su último valor observado. Población base y mortalidad: ONU, WPP 2024; México: INEGI, Encuesta Intercensal 2025. Detalle al final.}
\end{frame}

% ===== 6. Envejecer antes de ser ricos (N) =====
\begin{frame}{Envejecer antes de ser ricos}
  \vspace{0.5em}
  \begin{columns}[c,onlytextwidth]
    \begin{column}{0.45\textwidth}\small
      \begin{itemize}\setlength{\itemsep}{0.6em}
        \item Europa tardó \dest{55 años} en pasar de 10\% a 20\% de población de 65 años y más. Francia y Suecia tardaron todavía más.
        \item América Latina y el Caribe hará ese mismo recorrido en alrededor de \dest{27 años}, entre 2024 y 2051.
        \item La región envejece antes de alcanzar el ingreso con el que otros países construyeron sus estados de bienestar. Tendrá menos tiempo y menos recursos para adaptar sus sistemas de pensiones, salud y cuidados.
      \end{itemize}
    \end{column}
    \begin{column}{0.51\textwidth}\centering
      {\scriptsize\bfseries Años para pasar de 10\% a 20\% de personas de 65 y más\par}\vspace{0.5em}
      \begin{tikzpicture}[x=0.05cm,y=0.48cm,font=\scriptsize]
        \foreach \n/\d/\c/\i/\l in {Corea/17/barra/0/17,Japón/21/barra/1/21,América Latina/27/oxido/2/27,España/49/barra/3/49,Europa/55/barra/4/55,Francia/68/barra/5/68 o más,Suecia/70/barra/6/70 o más} {
          \fill[\c] (0,-\i-0.34) rectangle (\d,-\i+0.34);
          \node[anchor=east] at (-1,-\i) {\n};
          \node[anchor=west] at (\d,-\i) {\l};
        }
      \end{tikzpicture}
    \end{column}
  \end{columns}
  \fuente{Cálculo propio con ONU, WPP 2024 (después de 2023 es proyección); América Latina y el Caribe: proyección propia. Francia y Suecia ya pasaban de 10\% en 1950, inicio de la serie.}
\end{frame}

% ===== 7. También cambia la forma en que nos enfermamos (C7, ejemplo mexicano) =====
\begin{frame}{\mx También cambia la forma en que nos enfermamos}
  \vspace{1.2em}
  \begin{itemize}\setlength{\itemsep}{0.35em}
    \item En México, las enfermedades no transmisibles ya explican \dest{60.1\%} de los años de vida perdidos por muerte prematura o discapacidad.
    \item Diabetes, hipertensión, cáncer: son padecimientos crónicos que se atienden durante años.
    \item El IMSS calcula que su prima de equilibrio para atención médica tendría que subir \dest{34\%} para cubrir sus gastos.
  \end{itemize}
  \vspace{0.6em}
  Esta presión va a crecer con o sin reforma. Viene de la demografía y del perfil epidemiológico, y pasa lo mismo en el resto de la región.
  \fuente{CIEP (2026), \emph{Implicaciones del Paquete Económico 2027}, con base en IHME, OMS e IMSS.}
\end{frame}

% ===== 8. Transición (C9) =====
\begin{frame}[plain]
  \vfill
  \noindent\makebox[\textwidth][c]{\begin{beamercolorbox}[wd=0.72\textwidth,sep=14pt,center]{bandacolor}
    {\color{marino}\bfseries\Large Esto ya no depende de nosotros.\par\vspace{0.5em}Lo que sigue, sí.\par}
  \end{beamercolorbox}}
  \vfill
\end{frame}

% ===== 9. Avances: la expansión no contributiva (P) =====
\begin{frame}{Avances: la expansión no contributiva}
  \vspace{0.7em}
  \begin{itemize}\setlength{\itemsep}{0.4em}
    \item En este cuarto de siglo, el mayor avance en la protección de la vejez vino de las \clave{pensiones no contributivas}. Ampliaron la cobertura de manera real y llegaron a personas que nunca habrían reunido las cotizaciones de una pensión contributiva.
    \item Se hizo por dos caminos: como piso focalizado en algunos países y como beneficio universal en otros.
    \item En los dos casos el costo fiscal crece cada año, con el número de personas mayores. Y si la pensión no contributiva no se integra con la contributiva, puede debilitar el incentivo a cotizar.
  \end{itemize}
  \banda{El gran avance del cuarto de siglo fue no contributivo: se amplió la cobertura sin ampliar la base que la paga.}
  \vspace{-0.9em}
  \fuente{Villarreal (2026), «América Latina necesita una nueva seguridad social», XXXVIII Seminario Regional de Política Fiscal, CEPAL, Santiago.}
\end{frame}

% ===== 10. Transformaciones: el péndulo de las reformas (N) =====
\begin{frame}{Transformaciones: el péndulo de las reformas}
  \centering
  \vspace{0.3em}
  \begin{tikzpicture}[x=1.45cm,y=1cm,font=\scriptsize]
    \draw[pie,line width=0.8pt,-{Straight Barb[length=3pt]}] (-0.4,0) -- (8.75,0);
    \foreach \x/\a/\p/\t/\c in {
      0/1981/Chile/capitalización\\individual/marino,
      2/2008/Chile/pilar\\solidario/medio,
      4/2010/Bolivia/gestión\\pública/medio,
      6/desde 2020/Chile y Perú/retiros\\anticipados/medio,
      8/2025/Chile/seguro social y\\más cotización/oxido} {
      \fill[\c] (\x,0) circle (2.6pt);
      \draw[\c] (\x,0) -- (\x,0.45);
      \node[anchor=south,align=center,text width=2.3cm] at (\x,0.45) {{\color{\c}\bfseries\small \a}\\[0.1em]\textbf{\p}\\ \t};
    }
    \foreach \x/\a/\p/\t/\c in {
      1/1997/México/cuentas\\individuales/marino,
      3/2008/Argentina/regreso\\al reparto/medio,
      5/2020/México/más aportación\\patronal/medio,
      7/2024/Colombia/sistema\\de pilares/oxido} {
      \fill[\c] (\x,0) circle (2.6pt);
      \draw[\c] (\x,0) -- (\x,-0.45);
      \node[anchor=north,align=center,text width=2.3cm] at (\x,-0.45) {{\color{\c}\bfseries\small \a}\\[0.1em]\textbf{\p}\\ \t};
    }
  \end{tikzpicture}\par
  \vspace{0.1em}
  \notagraf{\centering Línea de tiempo sin escala.}
  \vspace{-0.3em}
  \banda{\small La región pasó de reformar parámetros, o de ir y venir entre modelos, a preguntarse por la arquitectura.}
  \vspace{-1.0em}
  \fuente{Leyes de reforma de cada país. Chile, 2025: Ley 21.735. Colombia, 2024: Ley 2381.}
\end{frame}

% ===== 11. Chile 2025 y Colombia 2024 (N) =====
\begin{frame}{Chile 2025 y Colombia 2024: dos respuestas al mismo problema}
  \vspace{0.1em}
  \begin{columns}[T,onlytextwidth]
    \begin{column}{0.49\textwidth}
      \caja{Chile: Ley 21.735 (marzo de 2025)}{%
        Nueva cotización del empleador de \textbf{7\%} de la renta imponible. Es gradual: un punto desde agosto de 2025 y completa en 9 años, extensible a 11.\par\vspace{0.3em}
        Con el 1.5\% del seguro de invalidez y sobrevivencia suma \textbf{8.5\%}: 4.5\% va al ahorro individual y 4\% a un nuevo Seguro Social.\par\vspace{0.3em}
        La Pensión Garantizada Universal sube a 250 mil pesos, de forma escalonada por edad, entre septiembre de 2025 y septiembre de 2027.}
    \end{column}
    \begin{column}{0.49\textwidth}
      \caja{Colombia: Ley 2381 de 2024}{%
        Crea un \textbf{sistema de pilares}.\par\vspace{0.3em}
        En agosto de 2026 la Corte Constitucional declaró exequible la mayor parte de la reforma y devolvió nueve artículos a la Cámara para subsanar vicios de trámite.\par\vspace{0.3em}
        La entrada en vigor está prevista para el \textbf{1 de abril de 2027}.}
    \end{column}
  \end{columns}
  \banda{Las dos combinan un componente solidario con el ahorro individual.}
  \vspace{-1.0em}
  \fuente{Chile: Ley 21.735; Ministerio de Hacienda de Chile, 2025. Colombia: Ley 2381 de 2024; Corte Constitucional de Colombia, 2026.}
\end{frame}

% ===== 12. Lo que sigue pendiente en pensiones (P) =====
\begin{frame}{Lo que sigue pendiente en pensiones}
  \vspace{0.5em}
  \begin{columns}[c,onlytextwidth]
    \begin{column}{0.47\textwidth}\small
      \begin{itemize}\setlength{\itemsep}{0.6em}
        \item \clave{Suficiencia.} Falta saber qué pensión tendrá la generación que se retira sólo con su cuenta individual.
        \item \clave{Cobertura contributiva.} En México, Colombia, Ecuador, Paraguay y República Dominicana cotiza o está afiliada menos de la mitad de los ocupados. Costa Rica y Uruguay son la excepción.
        \item \clave{Informalidad.} Cerca de la mitad del empleo de la región es informal: \dest{46.7\%} en 2025, con diferencias enormes entre países.
      \end{itemize}
      \vspace{0.3em}
      Las promesas de pensión maduraron antes de que se ampliara la base que cotiza.
    \end{column}
    \begin{column}{0.49\textwidth}\centering
      \begin{tikzpicture}
        \begin{axis}[base,width=0.92\linewidth,height=4.7cm,xbar,bar width=8pt,xmin=0,xmax=92,
            ymajorgrids=false,xmajorgrids,axis y line*=left,y axis line style={draw=none},
            ytick={1,...,7},ytick style={draw=none},
            yticklabels={Paraguay,México,Ecuador,Colombia,Rep.\ Dominicana,Costa Rica,Uruguay},
            ymin=0.4,ymax=7.6,xtick={0,25,50,75},xlabel={\% de ocupados que cotizan o están afiliados a pensiones}]
          \addplot[fill=oxido,draw=oxido,bar shift=0pt] coordinates {(25,1)(34,2)(35,3)(43,4)(43,5)};
          \addplot[fill=barra,draw=marino,bar shift=0pt] coordinates {(73,6)(77,7)};
          \node[font=\tiny,anchor=west] at (axis cs:25,1) {25};
          \node[font=\tiny,anchor=west] at (axis cs:34,2) {34};
          \node[font=\tiny,anchor=west] at (axis cs:35,3) {35};
          \node[font=\tiny,anchor=west] at (axis cs:43,4) {43};
          \node[font=\tiny,anchor=west] at (axis cs:43,5) {43};
          \node[font=\tiny,anchor=west] at (axis cs:73,6) {73};
          \node[font=\tiny,anchor=west] at (axis cs:77,7) {77};
        \end{axis}
      \end{tikzpicture}
    \end{column}
  \end{columns}
  \fuente{CEPALSTAT, ocupados que aportan o están afiliados a sistemas de pensiones, 2023--2024. OIT, \emph{Panorama Laboral 2025}: informalidad, primer semestre de 2025, promedio de 12 países. Villarreal (2026), CEPAL.}
\end{frame}

% ===== 13. Una arquitectura multipilar (K, P) =====
\begin{frame}{Una arquitectura multipilar}
  \centering
  \vspace{0.5em}
  \begin{tikzpicture}[font=\scriptsize,
      p/.style={draw=marino,line width=0.9pt,fill=banda,minimum width=2.75cm,minimum height=1.75cm,align=center,text width=2.5cm}]
    \node[p] (p0) at (0,0) {{\color{oxido}\bfseries\normalsize Pilar 0}\\[0.2em] No contributivo:\\ piso contra la pobreza\\ en la vejez};
    \node[p] (p1) at (3.5,0) {{\color{oxido}\bfseries\normalsize Pilar 1}\\[0.2em] Contributivo\\ obligatorio y solidario};
    \node[p] (p2) at (7.6,0) {{\color{marino}\bfseries\normalsize Pilar 2}\\[0.2em] Ahorro individual\\ obligatorio};
    \node[p] (p3) at (11.1,0) {{\color{marino}\bfseries\normalsize Pilar 3}\\[0.2em] Ahorro\\ voluntario};
    \draw[oxido,line width=1.2pt,{Straight Barb[length=3pt]}-{Straight Barb[length=3pt]}] (p0.east) -- (p1.west);
    \draw[marino,line width=1.2pt,{Straight Barb[length=3pt]}-{Straight Barb[length=3pt]}] (p2.east) -- (p3.west);
  \end{tikzpicture}\par
  \vspace{0.7em}
  {\color{marino}\bfseries\normalsize «El cero se habla con el uno, el dos con el tres.»\par}
  \vspace{0.7em}
  \raggedright
  \begin{itemize}\setlength{\itemsep}{0.3em}\small
    \item El piso no contributivo y la pensión contributiva se diseñan juntos, para que cada año cotizado sume a lo que la persona recibe.
    \item El ahorro obligatorio y el voluntario comparten cuentas, reglas e incentivos.
    \item Cuando cada pilar se reforma por separado, uno termina por debilitar al otro.
  \end{itemize}
\end{frame}

% ===== 14. ¿Cuánto queda para decidir? (C10) =====
\begin{frame}{\mx ¿Cuánto queda para decidir?}
  \vspace{0.1em}
  \begin{columns}[c,onlytextwidth]
    \begin{column}{0.47\textwidth}\small
      \begin{itemize}\setlength{\itemsep}{0.45em}
        \item En México, los ingresos disponibles del sector público son \dest{20.3\%} del PIB. Dejamos fuera el petróleo: lo que entra por el Fondo Mexicano del Petróleo (0.7) sale, en buena parte, como apoyo a Pemex (0.5).
        \item Les restamos lo que no se decide año con año: gasto no programable, pensiones y la operación propia de IMSS, ISSSTE y CFE.
        \item Queda \dest{4.0\%} del PIB, más deuda, para salud, educación, infraestructura, seguridad, cuidados y todo lo demás.
      \end{itemize}
      \vspace{0.3em}
      Son sólo sumas y restas, sin supuestos de comportamiento. Y 4.0\% es el máximo.
    \end{column}
    \begin{column}{0.50\textwidth}\centering
      \begin{tikzpicture}
        \begin{axis}[base,width=\linewidth,height=4.1cm,ybar,bar width=17pt,ymin=0,ymax=24,
            ytick={0,5,10,15,20},ylabel={\% del PIB},xmin=0.4,xmax=4.6,
            xtick={1,2,3,4},xtick style={draw=none},
            xticklabels={{Ingresos\\disponibles},{menos\\no programable},{menos\\pensiones},{menos operación\\de organismos}},
            x tick label style={font=\fontsize{4.6}{5.4}\selectfont,align=center}]
          \addplot[fill=barra,draw=marino,bar shift=0pt] coordinates {(1,20.3)(2,12.1)(3,6.0)};
          \addplot[fill=oxido,draw=oxido,bar shift=0pt] coordinates {(4,4.0)};
          \node[font=\tiny,anchor=south] at (axis cs:1,20.3) {20.3};
          \node[font=\tiny,anchor=south] at (axis cs:2,12.1) {12.1};
          \node[font=\tiny,anchor=south] at (axis cs:3,6.0) {6.0};
          \node[font=\tiny\bfseries,text=oxido,anchor=south] at (axis cs:4,4.0) {4.0};
        \end{axis}
      \end{tikzpicture}\par\vspace{0.2em}
      \notagraf{Restas: no programable 8.2 (costo financiero 3.8, participaciones 4.2, Adefas 0.2); pensiones 6.1; operación propia 2.1. El resultado sin redondear es 3.95; con las cifras redondeadas la resta da 3.9.}
    \end{column}
  \end{columns}
  \fuente{Villarreal, \emph{Contabilidad de la Muerte}, actualización 2026, con SHCP, Informe enero--septiembre 2026.}
\end{frame}

% ===== 15. La brecha en salud (C12) =====
\begin{frame}{\mx La brecha en salud}
  \vspace{0.1em}
  \begin{columns}[c,onlytextwidth]
    \begin{column}{0.47\textwidth}\small
      \begin{itemize}\setlength{\itemsep}{0.45em}
        \item En 2027 el gasto público en salud sería \dest{2.9\%} del PIB. Para llegar al 6\% que recomienda la OPS faltan \dest{3.1} puntos.
        \item Una reforma fiscal clásica completa (IVA, ISR y gastos fiscales) daría alrededor de \dest{1.5} puntos.
        \item La Miscelánea 2027 aporta alrededor de \dest{0.3} puntos.
      \end{itemize}
    \end{column}
    \begin{column}{0.50\textwidth}\centering
      \begin{tikzpicture}
        \begin{axis}[base,width=0.92\linewidth,height=3.5cm,xbar,bar width=9pt,xmin=0,xmax=3.7,
            ymajorgrids=false,xmajorgrids,axis y line*=left,y axis line style={draw=none},
            ytick={1,2,3},ytick style={draw=none},yticklabels={Miscelánea 2027,Reforma clásica,Brecha a 6\%},
            ymin=0.4,ymax=3.6,xtick={0,1,2,3},xlabel={puntos del PIB}]
          \addplot[fill=oxido,draw=oxido,bar shift=0pt] coordinates {(0.3,1)(1.46,2)(3.1,3)};
          \node[font=\tiny,anchor=west] at (axis cs:0.3,1) {$\approx$\,0.3};
          \node[font=\tiny,anchor=west] at (axis cs:1.46,2) {1.46};
          \node[font=\tiny,anchor=west] at (axis cs:3.1,3) {3.1};
        \end{axis}
      \end{tikzpicture}
    \end{column}
  \end{columns}
  \vspace{-0.2em}
  \banda{Una reforma fiscal ambiciosa, por sí sola, no alcanza para financiar la salud universal.}
  \vspace{-1.1em}
  \fuente{CIEP (2026); OPS (2014); estimación propia del paquete de reforma clásica. Miscelánea: estimación previa a las exclusiones aprobadas por el Senado; el rendimiento final es algo menor. Detalle de fuentes al final.}
\end{frame}

% ===== 16. Las pensiones le quitan espacio a la salud (C13) =====
\begin{frame}{\mx Las pensiones le quitan espacio a la salud}
  \vspace{0.3em}
  \begin{columns}[c,onlytextwidth]
    \begin{column}{0.53\textwidth}\small
      \begin{itemize}\setlength{\itemsep}{0.45em}
        \item El PEF 2027 asigna \dest{6.3\%} del PIB a pensiones y \dest{2.9\%} a salud: más de dos pesos a pensiones por cada peso a salud.
        \item Las pensiones contributivas crecen \dest{4.7\%} real, sobre todo por IMSS e ISSSTE. Es gasto comprometido desde hace años.
        \item Si no aumentan los ingresos, cada punto adicional en pensiones es un punto menos para salud.
      \end{itemize}
    \end{column}
    \begin{column}{0.40\textwidth}\centering
      \begin{tikzpicture}
        \begin{axis}[base,width=\linewidth,height=3.7cm,ybar,bar width=20pt,ymin=0,ymax=7.5,ymajorgrids=false,
            ytick={0,2,4,6},ylabel={\% del PIB},xmin=0.4,xmax=2.6,xtick={1,2},xtick style={draw=none},
            xticklabels={Pensiones,Salud}]
          \addplot[fill=oxido,draw=oxido,bar shift=0pt] coordinates {(1,6.3)};
          \addplot[fill=barra,draw=marino,bar shift=0pt] coordinates {(2,2.9)};
          \node[font=\tiny,anchor=south] at (axis cs:1,6.3) {6.3};
          \node[font=\tiny,anchor=south] at (axis cs:2,2.9) {2.9};
        \end{axis}
      \end{tikzpicture}
    \end{column}
  \end{columns}
  \vspace{-0.3em}
  \banda{\small En toda la región, cada punto adicional en pensiones compite con la salud.}
  \vspace{-1.1em}
  \fuente{\pef. CIEP (2026), \emph{Implicaciones del Paquete Económico 2027}, caps. 5 y 10: cifras del proyecto; el dictamen incluye reasignaciones menores.}
\end{frame}

% ===== 17. ¿Financiar todo con impuestos generales? (C14, ejemplo mexicano) =====
\begin{frame}{\mx ¿Financiar todo con impuestos generales?}
  \vspace{0.7em}
  \begin{itemize}\setlength{\itemsep}{0.35em}
    \item Un sistema financiado sólo con impuestos generales requiere una recaudación alta y estable.
    \item México recauda alrededor de \dest{18\%} del PIB, cerca de la mitad de lo que tenían los países que lograron cobertura universal por esa vía.
    \item Además, cada peso adicional que se recaude compite con las pensiones y con el costo de la deuda.
  \end{itemize}
  \banda{Con la recaudación actual, una salud universal financiada sólo con impuestos no es viable.}
  \vspace{0.2em}
  La pregunta es cómo ampliar el aseguramiento, no cómo sustituirlo.
\end{frame}

% ===== 18. El gasto de bolsillo en la región (C16, regionalizada) =====
\begin{frame}{Dos sistemas y el gasto de bolsillo}
  \vspace{0.5em}
  \begin{columns}[c,onlytextwidth]
    \begin{column}{0.46\textwidth}\small
      \begin{itemize}\setlength{\itemsep}{0.6em}
        \item En América Latina y el Caribe, \dest{29.7\%} del gasto en salud sale directamente del bolsillo de los hogares. En México es \dest{41.2\%}; en Colombia, \dest{14.6\%}.
        \item El gasto público en salud de la región es \dest{4.1\%} del PIB. Va de \dest{2.7\%} en México a \dest{6.5\%} en Uruguay.
        \item Donde conviven un sistema contributivo y otro para quien no cotiza, lo que ninguno de los dos cubre lo paga el hogar, sin planearlo y cuando menos puede.
      \end{itemize}
    \end{column}
    \begin{column}{0.50\textwidth}\centering
      \begin{tikzpicture}
        \begin{axis}[base,width=0.92\linewidth,height=5.0cm,xbar,bar width=7pt,xmin=0,xmax=50,
            ymajorgrids=false,xmajorgrids,axis y line*=left,y axis line style={draw=none},
            ytick={1,...,9},ytick style={draw=none},
            yticklabels={Colombia,Uruguay,Costa Rica,Argentina,Rep.\ Dominicana,América Latina y el Caribe,Ecuador,Paraguay,México},
            ymin=0.4,ymax=9.6,xtick={0,10,20,30,40},xlabel={gasto de bolsillo, \% del gasto corriente en salud, 2023}]
          \addplot[fill=barra,draw=marino,bar shift=0pt] coordinates {(14.6,1)(17.1,2)(24.1,3)(24.5,4)(25.1,5)(30.9,7)(36.3,8)(41.2,9)};
          \addplot[fill=oxido,draw=oxido,bar shift=0pt] coordinates {(29.7,6)};
          \foreach \v/\y in {14.6/1,17.1/2,24.1/3,24.5/4,25.1/5,29.7/6,30.9/7,36.3/8,41.2/9} {
            \edef\tmp{\noexpand\node[font=\noexpand\tiny,anchor=west] at (axis cs:\v,\y) {\v};}\tmp
          }
        \end{axis}
      \end{tikzpicture}
    \end{column}
  \end{columns}
  \fuente{OMS, Global Health Expenditure Database, y Banco Mundial, 2023: gasto de bolsillo y gasto en salud del gobierno general.}
\end{frame}

% ===== 19. El dilema Gruber (C17) =====
\begin{frame}{El dilema Gruber}
  \vspace{0.9em}
  \begin{itemize}\setlength{\itemsep}{0.35em}
    \item En un sistema mixto es muy difícil dar el mismo servicio de salud cuando unos contribuyen y otros no.
    \item Si el servicio es igual para todos, ¿para qué cotizar? Si es distinto, terminamos con dos sistemas.
    \item Un sistema de salud tiene que lograr dos cosas: que nadie se \clave{empobrezca} por enfermarse, y que \clave{convenga} estar en el sistema contributivo.
  \end{itemize}
  \vspace{0.6em}
  En México, como en el resto de la región, se ha avanzado en una o en otra, pero rara vez en las dos al mismo tiempo.
  \fuente{El planteamiento es de Jonathan Gruber (MIT).}
\end{frame}

% ===== 20. El costo de formalizarse (C18) =====
\begin{frame}{El costo de formalizarse}
  \vspace{0.5em}
  \begin{itemize}\setlength{\itemsep}{0.3em}
    \item Al formalizarse, un trabajador empieza a pagar cuotas e impuestos y, en algunos casos, deja de recibir apoyos.
    \item El Banco Mundial encuentra que, en América Latina, la mayor parte de ese costo corresponde a las \clave{contribuciones de seguridad social}.
    \item Quién paga realmente las cuotas, el trabajador o la empresa, depende de cómo se ajustan salarios y empleo. Para México no lo tenemos bien medido.
    \item Lo que sí sabemos es que no todas las cuotas pesan igual: la \clave{cuota fija} pesa proporcionalmente mucho más en los salarios bajos.
  \end{itemize}
  \vspace{0.5em}
  Si el trabajador no valora lo que recibe a cambio, la cuota se percibe como un impuesto más. El diseño puede cambiar eso.
  \fuente{Fietz et al. (2025), \emph{(In)Formalizing Jobs in LAC}, Banco Mundial. Detalle de fuentes al final.}
\end{frame}

% ===== 21. El escalón (C21) =====
\begin{frame}{El escalón}
  \vspace{0.5em}
  \begin{columns}[c,onlytextwidth]
    \begin{column}{0.50\textwidth}\small
      \begin{itemize}\setlength{\itemsep}{0.6em}
        \item Al pasar a un empleo formal, el ingreso neto puede caer: se empiezan a pagar cuotas y se pueden perder apoyos. Ese salto es la \clave{trampa de la formalización}.
        \item El escalón es más alto en los ingresos medios y bajos, donde la cuota fija pesa más y los apoyos importan más.
        \item Uruguay reconoció el desincentivo. Desde 2022 suspendió el control del tope de ingreso de sus asignaciones familiares, y entraron \dest{21 mil} hogares adicionales.
      \end{itemize}
    \end{column}
    \begin{column}{0.46\textwidth}\centering
      \begin{tikzpicture}[x=0.5cm,y=0.5cm,font=\tiny]
        \draw[-{Straight Barb[length=2pt]},pie] (0,0) -- (0,6.4);
        \draw[-{Straight Barb[length=2pt]},pie] (0,0) -- (10.6,0);
        \node[anchor=south] at (0.6,6.4) {ingreso neto};
        \node[anchor=north east] at (10.6,-0.1) {ingreso formal};
        \draw[marino,thick,dashed] (0,2) -- (10,5.6);
        \draw[oxido,thick] (0,2.6) -- (4.5,4.0);
        \draw[oxido,thick] (4.5,2.9) -- (10,4.9);
        \draw[oxido,thick,densely dotted] (4.5,4.0) -- (4.5,2.9);
        \node[text=oxido,anchor=west] at (4.7,2.3) {el escalón};
        \node[text=marino,anchor=south east] at (9.8,5.7) {sin saltos};
      \end{tikzpicture}\par\vspace{0.6em}
      \notagraf{\centering Esquema conceptual, sin escala.}
    \end{column}
  \end{columns}
  \fuente{Fietz et al. (2025), Banco Mundial; BPS Uruguay, resoluciones 2021--2026; MEF Uruguay, Rendición de Cuentas 2023. Detalle de fuentes al final.}
\end{frame}

% ===== 22. Bismarck, Beveridge y nuestra propuesta (C24) =====
\begin{frame}{Bismarck, Beveridge y nuestra propuesta}
  \vspace{-0.2em}
  {\scriptsize
  \begin{tabularx}{\textwidth}{@{}>{\raggedright\arraybackslash}p{2.1cm}>{\raggedright\arraybackslash}p{2.0cm}>{\raggedright\arraybackslash}p{1.6cm}>{\raggedright\arraybackslash}X@{}}
    \toprule
    \textbf{Decisión} & \textbf{Bismarck} & \textbf{Beveridge} & \textbf{Nuestra propuesta}\\
    \midrule
    ¿Con qué se financia? & Cuotas obrero-patronales & Impuestos generales & Se mantienen las cuotas obrero-patronales; quien no cotiza paga una prima según su ingreso; el gobierno completa\\\addlinespace
    ¿La prima es de cada persona? & Sí & No & Sí\\\addlinespace
    ¿Quién la paga? & Trabajador y patrón (en México, también el gobierno) & El gobierno & Asalariados, igual que hoy. Los demás, según su ingreso; el gobierno paga la prima completa sólo a los más pobres\\\addlinespace
    ¿El derecho depende del pago? & Sí & No, basta con residir & Sí, y todos tienen su prima cubierta\\
    \bottomrule
  \end{tabularx}}
  \vspace{-0.3em}
  \banda{Nuestra propuesta combina de otra manera las mismas cuatro decisiones.}
\end{frame}

% ===== 23. Cómo funcionaría (C25) =====
\begin{frame}{Cómo funcionaría}
  \vspace{0.5em}
  \begin{itemize}\setlength{\itemsep}{0.35em}
    \item Quien tiene empleo formal sigue pagando sus \clave{cuotas obrero-patronales}, como hoy. Lo que cambia es que debe ver con claridad qué paga y qué recibe.
    \item Quien no cotiza se inscribe en los seguros médicos del IMSS y paga una \clave{prima según su ingreso}. El gobierno completa la diferencia.
    \item Sólo a la población más pobre el gobierno le paga la prima completa, como un impuesto negativo. Si la pagara para todos, se perdería el incentivo a cotizar.
  \end{itemize}
  \vspace{-0.2em}
  \banda{Todos dentro del mismo sistema de aseguramiento, sin esquemas paralelos.}
  \vspace{-0.9em}
  \banda{\small La propuesta no sube las cuotas obrero-patronales. Con el cambio en Infonavit, la tasa legal sobre la nómina bajaría 3 puntos.}
\end{frame}

% ===== 24. Cómo decide un trabajador (C20; ocupa el lugar de «Qué mueve la decisión de formalizarse») =====
\begin{frame}{Qué mueve la decisión de formalizarse}
  \vspace{0.1em}
  \begin{columns}[T,onlytextwidth]
    \begin{column}{0.49\textwidth}
      \caja{Lo que cuesta ser formal}{%
        Cuotas de seguridad social: en América Latina son el componente mayor.\par\vspace{0.3em}
        Impuestos sobre el ingreso.\par\vspace{0.3em}
        Apoyos y servicios no contributivos que se pueden perder.}
    \end{column}
    \begin{column}{0.49\textwidth}
      \caja{Lo que recibe a cambio}{%
        Salud, pensión, vivienda, guarderías, incapacidades.\par\vspace{0.3em}
        Cuenta lo que el trabajador cree que vale, no lo que le cuesta al sistema.\par\vspace{0.3em}
        Si la cobertura sin cotizar se parece a la contributiva, cotizar vale menos.}
    \end{column}
  \end{columns}
  \vspace{0.5em}
  {\small
  \begin{itemize}\setlength{\itemsep}{0.25em}
    \item El costo neto de formalizarse es lo que se paga menos el valor que el trabajador asigna a lo que recibe. Para un asalariado de la región, va de \dest{20} a \dest{40\%} del ingreso del hogar.
    \item De ahí salen dos reglas de diseño: lo que se recibe debe depender de lo que se cotiza, y la cuota por el mismo ingreso no debe depender de si el trabajo es formal o informal.
    \item Todavía no hemos estimado cuánto responde la formalidad en México a estos cambios. Lo que sigue es lógica de diseño, no un resultado.
  \end{itemize}}
  \vspace{-0.5em}
  \fuente{Fietz et al. (2025), \emph{(In)Formalizing Jobs in LAC}, Banco Mundial, cap. 5: mediana en seis países; en México los impuestos pesan más que en el resto.}
\end{frame}

% ===== 25. Consolidar transferencias: un impuesto negativo (C27 y P, segundo eje) =====
\begin{frame}{Consolidar transferencias: un impuesto negativo}
  \vspace{0.2em}
  \begin{columns}[c,onlytextwidth]
    \begin{column}{0.47\textwidth}\small
      \begin{itemize}\setlength{\itemsep}{0.6em}
        \item Los programas con límites de ingreso generan un escalón: quien rebasa el límite pierde el apoyo completo.
        \item Con un \clave{impuesto negativo sobre la renta}, el apoyo baja poco a poco conforme aumenta el ingreso, y formalizarse siempre conviene.
        \item Esto requiere medir bien los ingresos. La factura electrónica y los pagos digitales ya lo hacen posible en México, Brasil, Uruguay y Colombia.
      \end{itemize}
    \end{column}
    \begin{column}{0.49\textwidth}\centering
      \begin{tikzpicture}[x=1cm,y=1cm,font=\tiny]
        \draw[-{Straight Barb[length=2pt]},pie] (0,0) -- (0,2.7);
        \draw[-{Straight Barb[length=2pt]},pie] (0,0) -- (4.9,0);
        \node[anchor=south] at (0,2.75) {beneficio};
        \node[anchor=west] at (4.95,0) {ingreso};
        \draw[oxido,thick] (0.2,2.15) -- (2.2,2.15);
        \draw[oxido,thick,dashed] (2.2,2.15) -- (2.2,0.3);
        \draw[oxido,thick] (2.2,0.3) -- (4.5,0.3);
        \draw[marino,thick] (0.2,1.85) -- (4.5,0.15);
        \node[text=oxido,anchor=south west] at (0.55,2.17) {sistema actual};
        \node[text=marino,anchor=south west] at (2.8,0.95) {impuesto negativo};
        \node[text=pie,anchor=north] at (2.2,-0.08) {umbral};
      \end{tikzpicture}\par\vspace{0.6em}
      \notagraf{\centering Esquema conceptual, sin escala.}
    \end{column}
  \end{columns}
  \vspace{-0.2em}
  \banda{\small El impuesto negativo es un puente hacia la contribución propia, no el destino.}
  \vspace{-1.1em}
  \fuente{Villarreal (2026), CEPAL.}
\end{frame}

% ===== 26. Nuevos instrumentos para financiar la salud (P, primer eje, y C30-31) =====
\begin{frame}{Nuevos instrumentos para financiar la salud}
  \vspace{0.1em}
  \begin{columns}[T,onlytextwidth]
    \begin{column}{0.32\textwidth}
      \caja{1. Cuotas rediseñadas}{%
        Canalizadas a salud, con un vínculo claro entre lo que se paga y lo que se recibe.\par\vspace{0.3em}
        Que el trabajador perciba un valor y no un impuesto.}
    \end{column}
    \begin{column}{0.32\textwidth}
      \caja{2. Primas según el ingreso}{%
        Para quien no cotiza: paga al menos una parte de su prima, proporcional a su ingreso.\par\vspace{0.3em}
        El gobierno paga la prima completa sólo a los más vulnerables.}
    \end{column}
    \begin{column}{0.32\textwidth}
      \caja{3. Seguros de ciclo de vida}{%
        Se pagan durante la vida laboral y cubren el gasto en salud de la vejez.\par\vspace{0.3em}
        Descargan al presupuesto público en el largo plazo.}
    \end{column}
  \end{columns}
  \vspace{0.5em}
  {\small
  \begin{itemize}\setlength{\itemsep}{0.25em}
    \item El gasto privado tiene que ser \clave{prepago} y no gasto de bolsillo. El prepago reparte el riesgo entre muchos y se paga antes de enfermarse.
    \item Caben también los seguros médicos privados de bajo costo y la atención primaria privada, si se regulan y se conectan con el resto del sistema.
  \end{itemize}}
  \vspace{-0.3em}
  \banda{\small ¿Qué reglas necesita el sector privado para participar sin fragmentar el sistema?}
  \vspace{-1.1em}
  \fuente{Villarreal (2026), CEPAL; Villarreal (2024), «¿Se puede reformar la seguridad social? Parte II», \emph{Arena Pública}.}
\end{frame}

% ===== 27. Cuidados: el siguiente instrumento (N) =====
\begin{frame}{Cuidados: el siguiente instrumento}
  \vspace{0.7em}
  \begin{itemize}\setlength{\itemsep}{0.4em}
    \item Entre 2025 y 2050, las personas de 80 años y más pasarán de \dest{13} a \dest{37 millones} en la región. A esa edad la dependencia deja de ser ocasional y el cuidado puede volverse permanente.
    \item Pensiones, salud y cuidados compiten por el mismo espacio presupuestario. Conviene analizarlos como un solo sistema.
    \item Hay instrumentos posibles: un seguro que se pague durante la vida laboral y cubra en la vejez, con un fondo solidario como base y con espacio para que lo público y lo privado se complementen.
  \end{itemize}
  \banda{¿Puede un seguro pagado durante la vida laboral financiar los cuidados de la vejez?}
  \vspace{-0.9em}
  \fuente{Población de 80 años y más: proyección propia (detalle al final). Villarreal (2026), CEPAL; Villarreal (2024), \emph{Arena Pública}.}
\end{frame}

% ===== 28. Dos preguntas para cualquier propuesta (C32) =====
\begin{frame}{Dos preguntas para cualquier propuesta}
  \vspace{1.6em}
  \begin{enumerate}\setlength{\itemsep}{0.35em}
    \item ¿Se puede pagar durante toda la transición demográfica sin que la deuda se vuelva insostenible?
    \item ¿Fortalece o debilita el incentivo a cotizar?
  \end{enumerate}
  \vspace{0.6em}
  Si una propuesta falla en alguna de las dos, tiene un problema de fondo. Esto aplica a lo que propongo y a cualquier alternativa. Para contestar la primera hay que hacer la cuenta completa: salud, pensiones y reforma fiscal juntas.
\end{frame}

% ===== 29. El nacimiento de CONOSAFI (C34, ampliada) =====
\begin{frame}{El nacimiento de CONOSAFI}
  \vspace{0.3em}
  {\small
  \begin{itemize}\setlength{\itemsep}{0.4em}
    \item \clave{Por qué nace.} En la región se ha discutido mucho cómo ampliar la cobertura, y mucho menos cómo pagarla en los próximos veinte o treinta años.
    \item \clave{Quién la funda.} Tres instituciones:
      \begin{itemize}
        \item Tecnológico de Monterrey --- coordinación, Héctor Villarreal
        \item Universidad de los Andes, Colombia --- Ramiro Guerrero
        \item Pontificia Universidad Católica de Chile --- Pablo Celhay
      \end{itemize}
    \item \clave{Cuándo.} Se anunció en el Seminario Regional de Política Fiscal de la CEPAL, en mayo de 2026. La primera reunión de trabajo fue en San José, Costa Rica, del 5 al 7 de octubre de 2026.
    \item \clave{Hacia dónde va.} Tiene un horizonte de al menos cinco años y tres trabajos en curso, uno por cada institución fundadora.
  \end{itemize}}
  \vspace{-0.2em}
  \fuente{CONOSAFI: Comunidad Latinoamericana de Conocimiento en Sistemas de Salud y Financiamiento.}
\end{frame}

% ===== 30. Una agenda para la RIUSS (C, preguntas abiertas, ampliada) =====
\begin{frame}{Una agenda para la RIUSS}
  \vspace{0.6em}
  Cinco preguntas abiertas:
  \vspace{0.3em}
  \begin{enumerate}\setlength{\itemsep}{0.3em}
    \item ¿Qué parte del gasto de bolsillo podría convertirse en prepago?
    \item ¿Con qué reglas el financiamiento privado complementa al público en lugar de debilitarlo?
    \item ¿Cuánto puede ahorrar la detección temprana cuando las enfermedades son crónicas?
    \item ¿Cuánto espacio dejan las pensiones para financiar la salud en los próximos veinte años?
    \item ¿Con qué instrumentos se van a financiar los cuidados de largo plazo?
  \end{enumerate}
  \banda{Buscamos investigadores y datos comparables en la región.}
\end{frame}

% ===== 31. Cierre (C35) =====
\begin{frame}[plain]
  \vfill
  \noindent\makebox[\textwidth][c]{\begin{beamercolorbox}[wd=0.78\textwidth,sep=10pt,center]{bandacolor}
    {\color{marino}\bfseries\normalsize ¿Podemos diseñar sistemas de protección social que la gente valore lo suficiente para querer participar, y que nuestros países puedan pagar durante su transición demográfica?\par}
  \end{beamercolorbox}}
  \begin{center}
  \vspace{1.6em}
  {\small Héctor Juan Villarreal Páez · hjvp@itesm.mx\\[0.2em]
  ITED · Escuela de Gobierno y Transformación Pública · Tecnológico de Monterrey\par}
  \end{center}
  \vfill
\end{frame}

% ===== Fuentes =====
\begin{frame}{Fuentes (1/2)}
  \vspace{0.1em}
  {\fontsize{6.1}{7.5}\selectfont\raggedright
  \textbf{Demografía de la región.} Proyección propia de la población de América Latina y el Caribe (49 países y territorios), 2025--2060, por componentes, en grupos quinquenales de edad. La fecundidad de cada país se mantiene en su último valor observado: registros nacionales de 2022 a 2025 para 18 países, reunidos por Fernández-Villaverde, J. (2026), \emph{A Recent Population History of Latin America and the Caribbean}; México, 1.23 hijos por mujer (INEGI, Encuesta Intercensal 2025); en los demás países, la estimación de Naciones Unidas para 2025. Población base de México: INEGI, Encuesta Intercensal 2025. Población base de los demás países, mortalidad, migración y patrón de fecundidad por edad: Naciones Unidas, DESA, \emph{World Population Prospects 2024}, variante media. Con la fecundidad de Naciones Unidas, el mismo cálculo da 729 millones en 2050 (publicado: 730) y el cruce entre mayores de 65 y menores de 15 hacia 2046. Fernández-Villaverde, J. (2023), \emph{The Demographic Future of Humanity: Facts and Consequences}.\par\vspace{0.5em}
  \textbf{Caída de la fecundidad.} Villarreal, H. (2026), «Cohort Replacement, Not Cascade: Union Composition and the Rapid Fertility Collapse in Latin America», documento de trabajo, Tecnológico de Monterrey: registros vitales de INEC (Costa Rica) y DANE (Colombia); DANE, Gran Encuesta Integrada de Hogares, 2008--2024.\par\vspace{0.5em}
  \textbf{Velocidad del envejecimiento.} Cálculo propio con \emph{World Population Prospects 2024}: años entre 10\% y 20\% de población de 65 años y más. Para América Latina y el Caribe, proyección propia.\par\vspace{0.5em}
  \textbf{Pensiones.} Villarreal, H. (2026), «América Latina necesita una nueva seguridad social», XXXVIII Seminario Regional de Política Fiscal, CEPAL, Santiago, mayo de 2026. Chile: Ley 21.735, publicada el 26 de marzo de 2025; Ministerio de Hacienda de Chile (2025). Colombia: Ley 2381 de 2024; Corte Constitucional de Colombia (2026). CEPALSTAT, ocupados que aportan o están afiliados a sistemas de pensiones, 2023--2024 (Argentina no está disponible en la misma serie). OIT, \emph{Panorama Laboral 2025 de América Latina y el Caribe}.\par\vspace{0.5em}
  \textbf{Gasto en salud y gasto de bolsillo.} OMS, Global Health Expenditure Database, y Banco Mundial, datos de 2023: gasto en salud del gobierno general como \% del PIB y gasto de bolsillo como \% del gasto corriente en salud.\par\vspace{0.5em}
  \textbf{México: carga de enfermedad, presupuesto 2027, salud y pensiones.} PEF 2027: dictamen de la Cámara de Diputados. CIEP (2026), \emph{Implicaciones del Paquete Económico 2027}, 11 de septiembre de 2026: gasto en salud, carga de enfermedad (IHME, 2021; OMS, 2024), prima de equilibrio del IMSS (IMSS, 2024), pensiones y Miscelánea; sus cifras corresponden al proyecto de presupuesto. Meta de 6\% del PIB: OPS (2014), \emph{Estrategia para el acceso universal a la salud y la cobertura universal de salud}.\par\vspace{0.5em}
  \textbf{México: espacio disponible.} Villarreal, \emph{Contabilidad de la Muerte}, actualización con el Informe de Finanzas Públicas y Deuda Pública enero--septiembre 2026 (SHCP) y los Criterios Generales de Política Económica 2027. PIB nominal de 36,500 mmdp; el supuesto favorece a la cuenta oficial.\par}
\end{frame}

\begin{frame}{Fuentes (2/2)}
  \vspace{0.2em}
  {\fontsize{6.6}{8.4}\selectfont\raggedright
  \textbf{Reforma fiscal clásica en México.} Estimación propia del rendimiento de un paquete de reforma clásica (IVA, ISR y gastos fiscales), alrededor de 1.5 puntos del PIB.\par\vspace{0.6em}
  \textbf{Trampa de la formalización.} Fietz, Joubert, Ñopo, Ocampo, Packard, Posadas y Rodriguez Chamussy (2025), \emph{(In)Formalizing Jobs in Latin America and the Caribbean: Taxes, Benefits, and Labor Market Incentives}, Banco Mundial, cap. 5. Uruguay: Banco de Previsión Social, resoluciones de directorio 2021--2026; Ministerio de Economía y Finanzas, Rendición de Cuentas 2023.\par\vspace{0.6em}
  \textbf{Modelo de aseguramiento, instrumentos y metas de financiamiento.} Villarreal, H. (2026), CEPAL. Villarreal, H. (2024), «¿Se puede reformar la seguridad social? Parte II», \emph{Arena Pública}, 10 de septiembre de 2024. Dilema de Gruber: Jonathan Gruber (MIT).\par\vspace{0.6em}
  \textbf{Cuota fija (apéndice).} Ley del Seguro Social, art. 106, fr. I. IMSS, \emph{Informe al Ejecutivo Federal y al Congreso de la Unión sobre la situación financiera y los riesgos del IMSS 2025--2026}: cuota fija patronal de 20.40\% de la UMA. INEGI, UMA 2026: 117.31 pesos diarios. CONASAMI, salario mínimo general 2026: 315.04 pesos diarios.\par\vspace{0.6em}
  \textbf{Infonavit (apéndice).} Ejercicio propio, preliminar, sobre tasas legales. No es una propuesta del IMSS ni del Gobierno Federal. Antecedente: Villarreal, H. (2022), «El Infonavit, ¿qué sigue?», \emph{Arena Pública}, 20 de diciembre de 2022.\par\vspace{0.6em}
  \textbf{Arquitectura institucional en México (apéndice).} Diario Oficial de la Federación: 29-XI-2019 (Insabi), 31-VIII-2022 (IMSS-Bienestar), 29-V-2023 (reforma a la Ley General de Salud).\par\vspace{0.6em}
  \textbf{CONOSAFI.} Comunidad Latinoamericana de Conocimiento en Sistemas de Salud y Financiamiento: nota de presentación (octubre de 2026) y documento de presentación de la iniciativa; reunión de trabajo en San José, Costa Rica, 5--7 de octubre de 2026.\par}
\end{frame}

% ===== A1. La cuota fija pesa más donde menos se gana (C19) =====
\begin{frame}{Apéndice · México: la cuota fija pesa más donde menos se gana}
  \vspace{0.5em}
  \begin{columns}[c,onlytextwidth]
    \begin{column}{0.50\textwidth}\small
      \begin{itemize}\setlength{\itemsep}{0.6em}
        \item En el seguro de enfermedades y maternidad, el patrón paga por cada trabajador una cuota fija de \dest{20.40\%} de la UMA, sin importar el salario.
        \item Para un trabajador de un salario mínimo, esa cuota equivale a \dest{7.6\%} del salario. Para uno de diez salarios mínimos, a \dest{0.8\%}.
        \item La parte de la cuota que distorsiona es la que no guarda relación con lo que el trabajador recibe, y se concentra en los salarios bajos.
      \end{itemize}
    \end{column}
    \begin{column}{0.46\textwidth}\centering
      \begin{tikzpicture}
        \begin{axis}[base,width=\linewidth,height=4.4cm,ybar,bar width=18pt,ymin=0,ymax=10,
            ytick={0,2,4,6,8},ylabel={\% del salario},xmin=0.4,xmax=4.6,
            xtick={1,2,3,4},xtick style={draw=none},xticklabels={1,2,5,10},xlabel={salarios mínimos}]
          \addplot[fill=oxido,draw=oxido,bar shift=0pt] coordinates {(1,7.60)};
          \addplot[fill=barra,draw=marino,bar shift=0pt] coordinates {(2,3.80)(3,1.52)(4,0.76)};
          \node[font=\tiny\bfseries,text=oxido,anchor=south] at (axis cs:1,7.60) {7.6};
          \node[font=\tiny,anchor=south] at (axis cs:2,3.80) {3.8};
          \node[font=\tiny,anchor=south] at (axis cs:3,1.52) {1.5};
          \node[font=\tiny,anchor=south] at (axis cs:4,0.76) {0.8};
        \end{axis}
      \end{tikzpicture}\par
      \notagraf{\centering Cuota fija patronal como \% del salario. Valores de 2026: UMA de 117.31 pesos diarios y salario mínimo general de 315.04.}
    \end{column}
  \end{columns}
  \fuente{Ley del Seguro Social, art. 106, fr. I; IMSS, \emph{Informe al Ejecutivo Federal y al Congreso de la Unión 2025--2026}; INEGI, UMA 2026; CONASAMI, salario mínimo general 2026. Cálculos propios.}
\end{frame}

% ===== A2. Infonavit: la propuesta (C23) =====
\begin{frame}{Apéndice · México, Infonavit: la propuesta}
  \vspace{0.1em}
  \begin{columns}[T,onlytextwidth]
    \begin{column}{0.49\textwidth}
      \caja{Qué se propone}{%
        Eliminar la aportación patronal de \textbf{5\%} del salario base de cotización al Infonavit.\par\vspace{0.3em}
        Destinar \textbf{2 puntos} a un seguro solidario de salud y cuidados en la vejez: un fondo común, no cuentas individuales. Primero cubre el déficit de los gastos médicos de pensionados; lo que sobra va a cuidados de largo plazo.\par\vspace{0.3em}
        Los otros \textbf{3 puntos} se dejarían de cobrar: la tasa legal sobre la nómina baja 3 puntos.}
    \end{column}
    \begin{column}{0.49\textwidth}
      \caja{El Instituto y los trabajadores}{%
        El Infonavit conserva su patrimonio y sus obligaciones como mutualidad.\par\vspace{0.3em}
        El trabajador deja de acumular 5\% de su salario en la subcuenta de vivienda y recibe a cambio un seguro.\par\vspace{0.3em}
        Quien no usa crédito tendría menos saldo para su pensión. Ese efecto no lo hemos cuantificado, y tampoco el efecto sobre la formalidad.}
    \end{column}
  \end{columns}
  \vspace{0.5em}
  {\small El seguro de gastos médicos de pensionados del IMSS ya es deficitario, y el envejecimiento presiona justo ahí.\par}
  \vspace{-0.4em}
  \fuente{Ejercicio propio, preliminar; no es una propuesta del IMSS ni del Gobierno Federal. Son tasas legales: no miden quién soporta la cuota. Antecedente: Villarreal (2022), «El Infonavit, ¿qué sigue?», \emph{Arena Pública}.}
\end{frame}

% ===== A3. El IMSS como sistema nacional de seguros de salud (C26) =====
\begin{frame}{Apéndice · México: el IMSS como sistema nacional de seguros de salud}
  \vspace{0.9em}
  \begin{itemize}\setlength{\itemsep}{0.3em}
    \item Partimos de lo que existe. El IMSS tiene los hospitales, el personal, buena parte de la formación de especialistas y el registro de derechohabientes. Para la población sin seguridad social, en cambio, el modelo cambió tres veces entre \dest{2019} y \dest{2023}.
    \item Que el IMSS agrupe los servicios de salud ligados al empleo formal: ISSSTE, Pemex e incluso fuerzas armadas.
    \item Que se dedique sólo a salud. Pensiones y cuidados, en otras instancias.
    \item Que pueda funcionar como caja de compensación y contratar servicios privados con recursos públicos.
  \end{itemize}
  \fuente{Diario Oficial de la Federación, 29-XI-2019, 31-VIII-2022 y 29-V-2023; Villarreal (2024). Detalle de fuentes al final.}
\end{frame}

% ===== A4. ¿De dónde saldría el 6 % del PIB? (C28) =====
\begin{frame}{Apéndice · México: ¿de dónde saldría el 6\% del PIB?}
  \vspace{0.1em}
  \begin{columns}[c,onlytextwidth]
    \begin{column}{0.48\textwidth}\small
      \begin{itemize}\setlength{\itemsep}{0.4em}
        \item \clave{Cuotas de salud:} de \dest{1.5\%} a \dest{2.5\%} del PIB, con el crecimiento de la masa salarial y la formalización.
        \item \clave{Gobierno federal}, sin contar cuotas: de \dest{1.4\%} a \dest{2.5\%}, sumando \dest{0.1} puntos cada año. La fuente serían los gastos fiscales.
        \item \clave{Estados y Ciudad de México:} \dest{1\%} del PIB, con participaciones e ingresos propios, sobre todo predial.
      \end{itemize}
      \vspace{0.3em}
      Nada de esto funciona si no se contiene el gasto en pensiones.
    \end{column}
    \begin{column}{0.48\textwidth}\centering
      \begin{tikzpicture}
        \begin{axis}[base,width=0.95\linewidth,height=4.2cm,ybar stacked,bar width=26pt,ymin=0,ymax=7.3,
            ytick={0,2,4,6},ylabel={\% del PIB},xmin=0.4,xmax=2.6,xtick={1,2},xtick style={draw=none},
            xticklabels={Hoy,Meta 2036},
            legend style={at={(0.5,-0.2)},anchor=north,legend columns=3},area legend]
          \addplot[fill=barra,draw=marino] coordinates {(1,1.5)(2,2.5)};
          \addlegendentry{Cuotas}
          \addplot[fill=medio,draw=marino] coordinates {(1,1.4)(2,2.5)};
          \addlegendentry{Federal}
          \addplot[fill=oxido,draw=oxido] coordinates {(1,0)(2,1.0)};
          \addlegendentry{Estados}
          \node[font=\tiny\bfseries,text=oxido,anchor=south] at (axis cs:1,2.9) {2.9};
          \node[font=\tiny\bfseries,text=oxido,anchor=south] at (axis cs:2,6.0) {6.0};
        \end{axis}
      \end{tikzpicture}\par
      \notagraf{\centering Hoy: aportación estatal propia no desglosada.}
    \end{column}
  \end{columns}
  \fuente{Villarreal (2024), «¿Se puede reformar la seguridad social? Parte II», \emph{Arena Pública}.}
\end{frame}

% ===== A5. Salud en el presupuesto 2027 (C15) =====
\begin{frame}{Apéndice · México: salud en el presupuesto 2027}
  \vspace{1.0em}
  \begin{itemize}\setlength{\itemsep}{0.35em}
    \item El gasto en salud aumenta \dest{10.6\%} real y llega a \dest{1.14 billones} de pesos. Si se descuenta la inflación de los servicios médicos, el aumento es de \dest{8.6\%}.
    \item El \dest{75.6\%} del aumento se va a población con seguridad social: IMSS crece \dest{11.6\%}, ISSSTE \dest{18.8\%} e IMSS-Bienestar \dest{9.0\%}.
    \item De ese total, \dest{760 mil millones} (\dest{67\%}) son para la población con seguridad social y \dest{376 mil millones} (\dest{33\%}) para el resto. Para el resto, el aumento es de \dest{7.6\%} real.
  \end{itemize}
  \vspace{0.6em}
  El aumento recupera parte de los recortes recientes, pero la estructura sigue igual.
  \fuente{\pef. CIEP (2026), \emph{Implicaciones del Paquete Económico 2027}, cap. 5, con base en SHCP: cifras del proyecto; el dictamen incluye reasignaciones menores.}
\end{frame}

\end{document}
